% !Mode:: "TeX:UTF-8"

%? ************************************************************* %?
%?                                                               %?
%?                   您将在此处完成中英文摘要的填写                   %?
%?                                                               %?
%? ************************************************************* %?

% ====================== 摘要 ======================

\chapter*{\texorpdfstring{\centering \heiti \sanhao \textbf{摘\quad 要}}{摘要}}
\addcontentsline{toc}{chapter}{\texorpdfstring{\sihao{摘\qquad 要}}{摘要}}

\linespread{1.5}\selectfont
\fangsong\xiaosi

\vspace{1em}

\xfabstract{

    随着电子商务和智能物流的快速发展，传统人工仓储模式在作业效率、响应速度和环境适应性方面已难以满足现代仓储系统的需求。为提高仓储搬运作业的自动化水平，本文围绕智能仓储场景下的举升式AGV小车展开研究，完成了总体方案设计、机械结构设计、电气系统设计、运动学建模及控制仿真分析。首先，结合仓储环境对通行性、灵活性和搬运效率的要求，确定了以双轮差速驱动为核心的举升式AGV总体方案。其次，完成了移动底盘、驱动轮、升降相关结构及罩壳等机械部分设计，并对支撑梁进行了强度校核和有限元分析，验证了结构方案的安全性与可行性。在电气系统方面，完成了基于STM32F103C8T6的主控系统、电机驱动模块、红外与超声传感器接口电路的选型与设计。随后，建立了差速驱动底盘的运动学模型，推导了机器人线速度、角速度与左右轮角速度之间的关系，并基于MATLAB对直线、圆弧、S形和8字等典型轨迹进行了仿真分析。最后，设计了基于PID的双轮差速控制器，并在扰动条件下进行了目标跟踪仿真验证。结果表明，所提出的AGV设计方案具有较好的结构合理性和控制可行性，可为后续样机开发与系统优化提供参考。

}

\noindent{\heiti\xiaosi\textbf{关键词：}} 
\fangsong{举升式AGV；智能仓储；差速驱动；运动学建模；PID控制}

% ====================== Abstract ======================

\chapter*{}
\vspace*{-1.15cm}
{\centering\bfseries\sanhao ABSTRACT\par}
\addcontentsline{toc}{chapter}{\sihao{ABSTRACT}}
\vspace{0.5cm}

\linespread{1.5}\selectfont
\rmfamily\xiaosi

\xfenabstract{
    With the rapid development of e-commerce and intelligent logistics, the traditional manual warehousing mode can no longer meet the requirements of modern storage systems in terms of operating efficiency, response speed, and environmental adaptability. To improve the automation level of warehouse handling operations, this thesis studies a lifting AGV for intelligent warehousing scenarios, and carries out overall scheme design, mechanical structure design, electrical system design, kinematic modeling, and control simulation analysis. First, according to the requirements of warehouse environments for maneuverability, accessibility, and handling efficiency, an overall scheme of a lifting AGV with dual-wheel differential drive is proposed. Then, the mechanical parts, including the mobile chassis, driving wheels, lifting-related structures, and shell, are designed, and the supporting beam is analyzed through strength checking and finite element analysis to verify the safety and feasibility of the structural scheme. In the electrical system, the main control system based on STM32F103C8T6, the motor drive module, and the interface circuits for infrared and ultrasonic sensors are selected and designed. After that, the kinematic model of the differential-drive chassis is established, and the relationships among the robot linear velocity, angular velocity, and the angular velocities of the left and right wheels are derived. Based on MATLAB, simulation analysis is carried out for typical trajectories such as straight line, circular arc, S-shape, and figure-eight paths. Finally, a dual-wheel differential PID controller is designed, and target tracking simulation is performed under disturbance conditions. The results show that the proposed AGV scheme has good structural rationality and control feasibility, which can provide a reference for subsequent prototype development and system optimization.
}

\vspace{1em}

\noindent\textbf{Key words:} lifting AGV; intelligent warehousing; differential drive; kinematic modeling; PID control

\setcounter{secnumdepth}{2}

